% -------------------------------------------------------
% File: MainFunction.tex
% Description: Main function of the assembled unit with
%              block diagram, wiring visuals, and power schematic
% -------------------------------------------------------

\chapter{Main Function}

The assembled unit implements real-time license plate acquisition on the NVIDIA Jetson Nano.  
The USB camera captures frames, the Jetson performs detection and optical character recognition,  
and results can be visualized on the HDMI display. This chapter describes the overall function,  
physical interconnects, and power path.

\section*{Overview}
\begin{itemize}
	\item \textbf{Compute Core:} Jetson Nano (4~GB) processes video frames locally.
	\item \textbf{Image Acquisition:} Microsoft LifeCam Show (USB) delivers the input stream.
	\item \textbf{User I/O:} HDMI monitor for display; USB keyboard and mouse for setup.
	\item \textbf{Storage:} \SI{32}{\giga\byte} microSD with JetPack~5.1.2.
	\item \textbf{Power:} \SI{5}{\volt} \SI{4}{\ampere} DC adapter.
\end{itemize}

% -------------------------------------------------------
% Table: Interfaces / Pins used
% -------------------------------------------------------
\begin{table}[H]
	\centering
	\caption{Interfaces Used by the Assembly}
	\begin{tabular}{|L{3.2cm}|L{4.2cm}|L{5.8cm}|}
		\hline
		\textbf{Subsystem} & \textbf{Jetson Port} & \textbf{Notes} \\ \hline
		USB Camera & USB~2.0/3.0 Type-A & Connect directly to a USB port on Jetson Nano. \\ \hline
		Display & HDMI & Use a short high-quality HDMI cable. \\ \hline
		User Input & USB Type-A & Keyboard and mouse on separate ports or via USB hub. \\ \hline
		Network (optional) & Ethernet (RJ45) & Recommended for first setup and package updates. \\ \hline
		Power & DC Barrel Jack & Official \SI{5}{\volt} \SI{4}{\ampere} adapter (5~V DC). \\ \hline
		Storage & microSD slot & \SI{32}{\giga\byte} microSD with JetPack~5.1.2. \\ \hline
	\end{tabular}
\end{table}

% -------------------------------------------------------
% TikZ 1: System Block Diagram
% -------------------------------------------------------
\begin{figure}[H]
	\centering
	\begin{tikzpicture}[
		every node/.style={font=\footnotesize, align=center},
		hw/.style={draw, rounded corners=3pt, fill=blue!10, minimum width=3.3cm, minimum height=1.2cm},
		proc/.style={draw, rounded corners=3pt, fill=orange!10, minimum width=3.3cm, minimum height=1.1cm},
		stor/.style={draw, rounded corners=3pt, fill=green!10, minimum width=3.3cm, minimum height=1.1cm},
		arr/.style={-{Latex[length=2mm]}, thick},
		>=Latex
		]
		
		% --- Hardware layer (bottom row) ---
		\node[hw] (camera) {USB Camera\\(Image Capture)};
		\node[hw, right=3.5cm of camera] (nano) {Jetson Nano\\(AI Inference Unit)};
		\node[hw, right=3.5cm of nano] (display) {HDMI Display\\(Visualization)};
		
		% --- Processing layer (above Nano) ---
		\node[proc, above=1.8cm of nano] (detect) {YOLOv8 Detection};
		\node[proc, above=1.8cm of detect] (ocr) {EasyOCR Recognition};
		\node[stor, right=3.5cm of ocr] (db) {SQLite Database\\(Local Logging)};
		
		% --- Arrows (data flow) ---
		\draw[arr] (camera) -- node[above]{Video Stream (USB)} (nano);
		\draw[arr] (nano) -- node[above]{HDMI Output} (display);
		
		\draw[arr] (nano.north) -- ++(0,0.4) -| (detect.south);
		\draw[arr] (detect) -- node[right]{Detected Region} (ocr);
		\draw[arr] (ocr) -- node[above]{Recognized Text} (db);
		
		% --- Power section ---
		\node[below=1.8cm of nano, font=\footnotesize] (pwr) {\SI{5}{\volt} \SI{4}{\ampere} DC Power Adapter};
		\draw[arr] (pwr) -- (nano.south);
		
		% --- Labels ---
		\node[below right=0.1cm and -1.1cm of detect, font=\scriptsize, align=left, text width=3.2cm]
		{On-device inference with TensorRT acceleration};
		
		\node[below right=-0.10cm and -0.5cm of ocr, font=\scriptsize, align=left, text width=3.5cm]
		{Character extraction and plate text decoding};
		
		\node[below right=0.1cm and -3.0cm of db, font=\scriptsize, align=left, text width=3.5cm]
		{Timestamps, detected IDs, billing logs stored locally};
		
	\end{tikzpicture}
	\caption{Overall license plate detection workflow implemented on Jetson Nano: from image capture to OCR and local database logging.}
\end{figure}

This diagram visually represents the complete end-to-end workflow of the system and the logical interaction between its key modules.  
It begins with the \textbf{USB camera}, which continuously captures live video frames from the environment. These frames are transmitted through the USB interface to the \textbf{Jetson Nano}, which serves as the system's central AI processing unit. Inside the Nano, the incoming image stream is processed by a pre-trained \textbf{YOLOv8 model} optimized for edge inference using TensorRT acceleration.  
Once a vehicle license plate is detected, the corresponding image region is forwarded to the \textbf{EasyOCR recognition layer}, which extracts alphanumeric text with high precision even under varying lighting conditions.  
The recognized plate numbers, along with timestamps, detection confidence scores, and event details, are then stored in the local \textbf{SQLite database} for persistent on-device logging and later retrieval.  
The output from the Jetson Nano is displayed on an \textbf{HDMI monitor} for live visualization of detection results, providing both verification and debugging feedback for the operator.  
The entire pipeline is powered by a regulated \SI{5}{\volt}, \SI{4}{\ampere} DC adapter that ensures stable GPU performance, thermal efficiency, and uninterrupted operation during continuous video inference.  
This workflow clearly outlines the system’s real-time nature—combining data acquisition, intelligent detection, optical recognition, and data management—all within a compact and power-efficient embedded platform.

% -------------------------------------------------------
% TikZ 2: Jetson Nano Top-View Connection Diagram (No Ethernet)
% -------------------------------------------------------
\begin{figure}[H]
	\centering
	\begin{tikzpicture}[
		every node/.style={font=\scriptsize, align=center},
		port/.style={draw, thick, rounded corners=2pt, fill=blue!10, minimum width=1.6cm, minimum height=0.7cm},
		gpio/.style={draw, thick, fill=orange!10, rounded corners=1pt, minimum width=0.45cm, minimum height=0.22cm},
		label/.style={font=\scriptsize},
		arr/.style={-{Latex[length=2mm]}, thick},
		node distance=0.6cm and 0.6cm
		]
		
		% --- Board outline ---
		\draw[thick, rounded corners=6pt, fill=gray!5] (0,0) rectangle (8.5,5);
		
		% --- Major Ports (spaced evenly) ---
		\node[port, fill=green!10] (usb1) at (1.2,4.3) {USB 3.0};
		\node[port, fill=green!10, right=0.6cm of usb1] (usb2) {USB 3.0};
		\node[port, fill=purple!10] (hdmi) at (7.4,3.4) {HDMI};
		\node[port, fill=yellow!30] (power) at (1.0,1.2) {DC Jack 5V};
		\node[port, fill=red!10] (cam) at (6.4,4.3) {CSI Camera Port};
		
		% --- GPIO Header (40-pin) ---
		\foreach \i in {0,...,19} {
			\node[gpio] (g\i) at (2.7+0.27*\i,0.5) {};
		}
		\node[label, below=0.3cm of g10, font=\scriptsize] {40-Pin GPIO Header};
		
		% --- microSD Slot ---
		\node[port, fill=yellow!15, rotate=90] (sd) at (0.6,2.8) {microSD Slot};
		
		% --- Arrows & Connection Labels ---
		\draw[arr] (cam.north) -- ++(0,0.9) node[above]{Camera Ribbon Cable};
		\draw[arr] (usb2.north) -- ++(0,1.0) node[above]{Keyboard / Mouse};
		\draw[arr] (usb1.north) -- ++(-0.2,1.0) node[above]{USB Camera};
		\draw[arr] (power.west) -- ++(-1.1,0) node[left]{5V 4A Power Adapter};
		\draw[arr] (hdmi.east) -- ++(1.2,0.3) node[right]{HDMI Monitor};
		\draw[arr, dashed] (g15.south) -- ++(0,-0.8) node[below]{Optional Sensor / Relay I/O};
		
		% --- Note ---
		\node[font=\tiny, align=left, text width=4.2cm, anchor=west] at (0.1,-1.15)
		{Note: GPIO header operates at 3.3V logic only.\\Avoid connecting 5V signals directly.};
		
	\end{tikzpicture}
	\caption{Top-view Jetson Nano hardware layout showing only the active connections used in the license plate detection system — USB camera, HDMI display, keyboard/mouse, CSI camera, GPIO, microSD, and 5V DC power.}
\end{figure}

The above top-view layout illustrates the exact hardware configuration of the assembled system.  
It shows how the USB camera, HDMI monitor, keyboard, and mouse are connected to the Jetson Nano through their respective ports, while power is supplied via a 5 V DC adapter.  
The CSI camera interface and GPIO header are highlighted for potential future extensions, and the microSD card slot provides storage for the operating system and project files.  
Arrows indicate the direction of power and data flow, ensuring that the physical connections match the functional workflow shown in the previous diagram.

% -------------------------------------------------------
% Table: Cable / Connector Checklist
% -------------------------------------------------------
\begin{table}[H]
	\centering
	\caption{Cable and Connector Checklist}
	\begin{tabular}{|L{4.2cm}|L{8.8cm}|}
		\hline
		\textbf{Connection} & \textbf{Checkpoints} \\ \hline
		USB Camera $\rightarrow$ Jetson & Firm seating, no strain; short cable preferred. \\ \hline
		HDMI $\rightarrow$ Monitor & Cable fully inserted; monitor set to correct input source. \\ \hline
		Keyboard / Mouse $\rightarrow$ Jetson & Separate ports or powered hub; verify LEDs (if present). \\ \hline
		Ethernet $\rightarrow$ Router/Switch & Link LEDs active; stable connection for first boot. \\ \hline
		DC Adapter $\rightarrow$ Jetson & Official \SI{5}{\volt} \SI{4}{\ampere}; avoid loose barrel jack. \\ \hline
	\end{tabular}
\end{table}